\documentclass[10pt,a4paper]{amsart}
\usepackage[margin=19mm]{geometry}
\usepackage{amsmath,amssymb,mathtools}
\usepackage[scr=rsfs,cal=euler]{mathalfa}
\usepackage{xfrac}
\usepackage[unicode,hidelinks,bookmarks=false]{hyperref}
\newcommand{\ie}{i.e.\ }
\newcommand{\cf}{cf.\ }
\newcommand{\resp}{respectively\ }
\newcommand{\Subsection}{\subsection}
\providecommand{\nobreakdash}{\nobreak\hbox{-}}
\setlength{\emergencystretch}{2em}
\multlinegap0pt
\usepackage{bm}
\usepackage{bbm}
\usepackage{dsfont}
\usepackage{xspace}
\usepackage{cancel}
\usepackage{mathtools}
\usepackage{amsthm}
\makeatletter
\@ifundefined{theorem}{\newtheorem{theorem}{Theorem}}{}
\@ifundefined{definition}{\newtheorem{definition}[theorem]{Definition}}{}
\@ifundefined{lemma}{\newtheorem{lemma}[theorem]{Lemma}}{}
\makeatother
\newcommand{\CalM}{\mathcal{M}}
\newcommand{\CalT}{\mathcal{T}}
\newcommand{\Eps}{\epsilon}
\newcommand{\Ni}{\noindent}
\newcommand{\R}{\mathbb R}
\newcommand{\T}{\mathbb{T}}
\newcommand{\abs}[1]{\left\vert#1\right\vert}
\newcommand{\dT}{\, {\rm d} T}
\newcommand{\del}{\partial}
\newcommand{\deta}{\, {\rm d} \eta}
\newcommand{\dtau}{\, {\rm d} \tau}
\newcommand{\dt}{\, {\rm d} t}
\newcommand{\dv}{\, {\rm d} v}
\newcommand{\dxi}{\, {\rm d} \xi}
\newcommand{\dx}{\, {\rm d} x}
\newcommand{\nn}{\nonumber}
\newcommand{\vint}[1]{\left\langle#1\right\rangle}
\newcommand{\vpran}[1]{\left(#1\right)}
\title[The $\mathcal{M}$-Operator and Uniqueness of Nonlinear Kinet]{The $\mathcal{M}$-Operator and Uniqueness of Nonlinear Kinetic Equations}
\author{Ricardo Alonso, Maria Pia Gualdani,, Weiran Sun}
\date{Source: arXiv:2506.20775v2 (CC BY 4.0)}
\begin{document}
\maketitle

\noindent\emph{Source and licence.} The $\mathcal{M}$-Operator and Uniqueness of Nonlinear Kinetic Equations. Version arXiv:2506.20775v2 (CC BY 4.0), distributed under the Creative Commons Attribution 4.0 International licence, as recorded in the arXiv licence metadata for this exact version and shown on the abstract page https://arxiv.org/abs/2506.20775v2. Full rights notice: licenses/arxiv-2506.20775-CC-BY-4.0.txt.\par\smallskip

\noindent\textit{Editorial scope.} Counted unit: the Theorem whose source label is \texttt{Thm\_Landau\_0} in the article (source0.tex lines 1788--1790, source label \texttt{Thm\_Landau\_0}), delivered together with the article's own complete original proof of it (source0.tex lines 1856--2287, one \texttt{proof} environment of 432 lines). No mathematics is written, repaired or re-derived: statement and proof are the article's own text line for line. The proof is detached from the statement in the source: the article states the result at lines 1788--1790 and prints its proof at lines 1856--2287, and the proof environment's own header names the statement, so the association is the article's own. Required context retained uncounted: trasp\_comm (lines 349--352); int\_xi\_M2 (lines 357--363); lapl (lines 357--363); lem:M-comm-x-v (lines 562--586); lem:kernels (lines 750--772); FPL (lines 1759--1762); D\_existence (lines 1775--1777); A\_sup (lines 1830--1834); nablaa\_L6 (lines 1830--1834). Every definition, display and label of those lines is retained unchanged. Omitted from this package and not redistributed: 1--348, 353--356, 364--561, 587--749, 773--1758, 1763--1774, 1778--1787, 1791--1829, 1835--1855, 2288--2441. Nothing inside a retained excerpt is omitted or altered: every retained body is delivered exactly as the article's lines are written, so no omission receipt of any kind accompanies this package. Citations: the retained excerpts carry no citation command at all, so no bibliography entry of the article is transcribed and no reference is redistributed. Reference adaptation: none was needed and none was made; every reference inside the delivered unit resolves inside it, so no unresolved reference or citation is produced. Printed slips kept literal: no printed word, symbol, hypothesis or number of the counted statement or of its proof was altered. Layout and class: the article's own class is \texttt{amsart} with packages including amsmath, amsfonts, amssymb, amscd, amsthm, amsbsy, bbm, epsf, calc, graphicx, esint, mathtools, color, xcolor; the wrapper is the lane's pinned \texttt{amsart} editorial wrapper at 10pt on A4, which restates the theorem-like environments the retained text needs and adds the article's own macro definitions that the retained text needs (\texttt{\textbackslash CalM}, \texttt{\textbackslash CalT}, \texttt{\textbackslash Eps}, \texttt{\textbackslash Ni}, \texttt{\textbackslash R}, \texttt{\textbackslash T}, \texttt{\textbackslash abs}, \texttt{\textbackslash dT}, \texttt{\textbackslash del}, \texttt{\textbackslash deta}, \texttt{\textbackslash dt}, \texttt{\textbackslash dtau}, \texttt{\textbackslash dv}, \texttt{\textbackslash dx}, \texttt{\textbackslash dxi}, \texttt{\textbackslash nn}, \texttt{\textbackslash vint}, \texttt{\textbackslash vpran}), with extra wrapper packages (bm, bbm, dsfont, xspace, cancel, mathtools). The delivered numbering is the unit's own. Assets: no retained span contains a figure environment, a graphics-inclusion command, a PDF-page inclusion, a table or a TikZ picture; no image file is copied into this package, the package declares no asset dependency and no image file is redistributed with it. Licence: version arXiv:2506.20775v2 of this article is distributed under the Creative Commons Attribution 4.0 International licence, as recorded in the arXiv licence metadata for this exact version and shown on the abstract page https://arxiv.org/abs/2506.20775v2. A full rights notice is delivered at licenses/arxiv-2506.20775-CC-BY-4.0.txt. Characters: every retained excerpt is pure ASCII, so the delivered engine, XeLaTeX with its default Latin Modern text font, reports no missing character.
\begin{align}
{\CalM} (\partial_t  h   -v \cdot \nabla_x   h)  
& = \vpran{\partial_t - v \cdot \nabla_x } ({\CalM} h) + [\CalM, \partial_t - v \cdot \nabla_x ]h,\label{trasp_comm}
\end{align}

\begin{align}
 & \int_t^{T_0}  \langle \xi \rangle ^2 {M}^2(t,T, \eta,\xi)  \dT 
\le    
 \frac{2}{\Eps \delta}, \label{int_xi_M2}
\\
 \int_0^{T_0}  \int_0^{T}  \int_{\R^3}  &|\xi|^2 | {M} \hat{h}| ^2d\xi  \dt \dT \le \frac{2}{\Eps \delta} \int_0^{T_0} \int_{\R^3} | \hat{h} |^2  \dxi \dt.  \label{lapl}
\end{align}

\begin{lemma} \label{lem:M-comm-x-v}
Suppose $h \in C^\infty_b(\T^3 \times \R^3)$. Then there exists an $L^2$-bounded operator $\CalT$ whose symbol belong
to $S^{-1}_{0,0}$ uniformly with respect to $\delta, \Eps, T-t$
such that 
\begin{align*}
 [\CalM, h] = \CalT \CalM,
\qquad
\CalT = a(x, v, D_x, D_v) \vint{D_v}^{-1},
\end{align*}
% with
% \begin{align*}
%  \CalT_1 = a_1(x, v, D_x, D_v) \vint{D_v}^{-1},
% \qquad 
%  \CalT_2 = T a_2 (x, v, D_x, D_v) \vint{D_x}^{-1},
% \end{align*}
where $T$ is the time in $\CalM$ and $a$ is a bounded operator in $L^2_{x,v}$ with bounds independent of $\delta, \epsilon, T, t$.
% \begin{align*}
% \norm{\CalT_1 \nabla_v g}_{L^2_{x, v}} 
% \leq C \|g\|_{L^2_{x, v}}, \qquad
% \norm{\CalT_2 g}_{L^2_{x, v}} 
% \leq C T \|g\|_{L^2_{x, v}}, \qquad  
% \forall g \in \mathcal{S}(\R^3).
% \end{align*}
% Here $C >0$ is independent of $\eta, \delta, \Eps, T-t$.
\end{lemma}

\begin{lemma} \label{lem:kernels}
Suppose $h \in C([0, T_0] \times \T^3)$ and $g \in C([0, T_0] \times \T^3 \times \R^3) \cap L^\infty([0, T_0] \times \T^3; L^\infty_{k_0}(\R^3))$ with $k_0 > 0$. 

\Ni (a) Suppose $h \geq 0$. Then there exists $\phi_a(x) \geq 0$ such that 
\begin{align*}
 h_a(t, x) := h * \phi_a \geq 0, 
\qquad
 h_a \to h \,\, \text{in} \,\,
C([0, T_0] \times \T^3)
\,\,\,
\text{as $a \to 0$}.
\end{align*}

\Ni (b) Let $\psi_a \in C^\infty_c(\T^3 \times \R^3)$ be a mollifier in $(x, v)$. Then 
\begin{align*}
 g_a (t, x, v)
:= g \ast \psi_a 
\to g \,\,\,
\text{in} \,\,\,
C([0, T_0] \times T^3 \times \R^3) \,\,\,
\text{as $a \to 0$.}
\end{align*}
\end{lemma}

\begin{align}\label{FPL}
\partial_t f + v \cdot \nabla_x f = \nu \Delta_v f  + Q_L ( f,f),  
% \textrm{div}_v \left( A[f] \nabla f - f \nabla a[f]\right)
\end{align}

\begin{definition} \label{D_existence}
A positive function $f$ is a classical solution to (\ref{FPL}) in $(0,T) \times \T^3 \times \R^{3}$ with initial data $f_{in} \in L^{\infty}_x L^\infty_{v,k}$ and $k>\frac{15}{2}$ if $f$ satisfies (\ref{FPL}) classically, $ \|\vint{v}^k f\|_{L^\infty_{t,x,v}}$ depends only on the initial data, and $ f \in C_{kin}^{2,\alpha}(\Omega)$ for any compact subset $\Omega \subset (0, T) \times \T^3 \times \mathbb{R}^{3}$.  
\end{definition} 

\begin{align}\label{A_sup}
\| A[f]\|_{L^\infty_v(\mathbb{R}^3)} & \le \| a[f]\|_{L^\infty_v(\mathbb{R}^3)}  \le c(m)  \|   \langle v \rangle^{m}f\|_{L^2_v(\mathbb{R}^3)} \quad \textrm{for some $m>3/2$, } \\
\| \nabla a [f]\|_{L^\infty_v(\mathbb{R}^3)}&  \le c(m) \|   \langle v \rangle^{m}f\|_{L^4_v(\mathbb{R}^3)} \quad \textrm{for some $m>3/4$, } \label{nablaA_sup}\\
\| \nabla a [f]\|_{L^6_v(\mathbb{R}^3)} & \le   \| f\|_{L^2_v(\mathbb{R}^3)}.\label{nablaa_L6} 
\end{align}

\begin{theorem} \label{Thm_Landau_0}
Let $f(x,v,t)$ be a solution of (\ref{FPL}) as in Definition \ref{D_existence} in the time interval $[0,T_\ast]$. Let $\nu >0$ be any arbitrary positive constant. Then $f$ is unique in the class of continuous functions $f$ such that  $ \langle v \rangle ^m f \in  L^\infty_{t,x}(L^4_v \cap L^1_v)$, for  $m > 3$. 
\end{theorem}

\begin{proof} (Proof of Theorem \ref{Thm_Landau_0}.)
Let $f_1$ and $f_2$ be two smooth solutions to (\ref{FPL}) with the same initial data $f_{in}$.  The function 
$$
w := g_1-g_2, \quad g_i = \langle v \rangle ^m f_i,
\qquad i = 1, 2,
$$ 
solves
\begin{align*}%\label{FPLg_U}
\partial_t w + v \cdot \nabla_x w = \;&\textrm{div}_v \left( A[f_1] \nabla w - w \nabla a[f_1]\right)+\textrm{div}_v \left( A[f_1-f_2] \nabla g_2 - g_2 \nabla a[f_1-f_2]\right)  +\nu \Delta_v w  \nonumber \\
&-m\; \textrm{div}_v \left( A[f_1]w \frac{v}{ \langle v \rangle ^2}\right)-m\; \textrm{div}_v \left( A[f_1-f_2]\; g_2 \frac{v}{ \langle v \rangle ^2}\right) \nonumber \\
&- m \left  \langle A[f_1]   \nabla w ,  \frac{v}{ \langle v \rangle ^2}\right  \rangle - m \left  \langle A[f_1-f_2]   \nabla g_2 ,  \frac{v}{ \langle v \rangle ^2}\right  \rangle \nonumber \\
& +m^2 w \left \langle A[f_1]  \frac{v}{ \langle v \rangle ^2}  ,  \frac{v}{ \langle v \rangle ^2}\right  \rangle +m^2 g_2 \left \langle A[f_1-f_2]  \frac{v}{ \langle v \rangle ^2}  ,  \frac{v}{ \langle v \rangle ^2}\right  \rangle  \nonumber \\
& + m w \left \langle  \nabla a[f_1],  \frac{v}{ \langle v \rangle ^2}\right  \rangle  + m \left \langle g_2 \nabla a[f_1-f_2],  \frac{v}{ \langle v \rangle ^2}\right  \rangle \nonumber  \\
&- 2m \nu  \;\frac{v \cdot \nabla w }{\langle v \rangle ^2} +\nu  (m^2+2m)\; \frac{|v|^2w }{\langle v \rangle ^4} -3\nu m \; \frac{ w }{\langle v \rangle ^2} \nonumber\\
%& +  \langle v \rangle ^m  Q_L ( h_1-h_2,M) + \langle v \rangle ^m  Q_L ( M,h_1-h_2) \nonumber\\
=:&\; I_1 + ...+ I_{14}.
\end{align*}

Apply the $\CalM$-operator to the above equation and obtain 
\begin{align*}
  \del_t \CalM w  + v \cdot \nabla_x \CalM w
= [\CalM, \partial_t + v \cdot \nabla_x]w + \sum_{i=1}^{14} \CalM  I_i   ,
\end{align*}
with the transport commutator defined in~\eqref{trasp_comm}
%\begin{align} \label{def:R-1}
%\widehat{R_1} = \CalF\vpran{[\CalM_n, \del_t + v \cdot \nabla_x] w_n}
%= \delta 2^{\beta n} \tfrac{\langle \xi \rangle^{2}}{1 + \delta 2^{\beta n} \int_t ^T \langle \xi +(t- \tau) \eta\rangle^{2} \;d\tau} \CalF\vpran{\CalM_n w_n}.
%\end{align}
%We apply the Fourier transform of the previous equation and 
We multiply the resulting equation by ${{\CalM} w}$ and integrate in $(0, T_0) \times (0,T) \times \T^3 \times \R^3$. This gives
\begin{align} \label{first_est}
\frac{1}{2}\int_0^{T_0} \int_0^{T}   \partial _t \iint_{\T^3 \times \R^3} |{\CalM} w|^2 
%\dx \dv \dt\dT
= & \int_0^{T_0}  \int_0^{T}  {\CalM} w  \iint_{\T^3 \times \R^3} [\CalM, \partial_t + v \cdot \nabla_x]w %\dx \dv \dt\dT 
\nonumber 
\\
&\;+ \int_0^{T_0}  \int_0^{T}  \iint_{\T^3 \times \R^3}  {{\CalM}w} {\CalM} ( {I_1} + ...+  {I_{14}}). %\dx \dv \dt\dT.
\end{align}


Since ${M}(\cdot,\cdot,T)=1$ and $\widehat w(\cdot,\cdot,0)=0$, the left hand side of (\ref{first_est}) reduces to 
$$
\frac{1}{2} \int_0^{T_0} \|\widehat w(T) \|_{L^2}^2 \dT = \frac{1}{2} \int_0^{T_0} \iint_{\T^3 \times \R^3}  w^2 \dx\dv\dt .
$$
The first term on the right hand side of (\ref{first_est}) can be estimated as 
\begin{align*}
 \delta \left(\frac{1}{2}+ \Eps\right)& \sum_{\eta}\int_0^{T_0} \int_0^T \int_{\R^3} \frac{ \langle \xi \rangle^{2}}{{ 1 + \delta \int_t ^T \langle \xi +(t- \tau) \eta\rangle^{2} \dtau}} |{M}\widehat w|^2 %\; d\xi  dtdT 
\\
 \le  & \;\delta  \left(\frac{1}{2}+ \Eps\right)\sum_{\eta} \left( \int_0^{T_0} \int_0^T \int_{\R^3}   
 | \xi |^{2}  |{M}\widehat w|^2 %\; d\xi dtdT 
 + T_0 \int_0^{T_0} \int_{\R^3}  |\widehat w|^2 
 %\; d\xi dt
 \right)
\\
 = & \;\delta  \left(\frac{1}{2}+ \Eps\right) \int_0^{T_0} \int_0^T \iint_{\T^3 \times \R^3} |\nabla_v \;\CalM w|^2 %dxdvdtdT 
 + T_0 \int_0^{T_0} \iint_{\T^3 \times \R^3}   w^2. %\;dxdvdt.
\end{align*}
Moreover, 
$$
 \int_0^{T_0} \int_0^T   \iint_{\T^3 \times \R^3} \CalM w  \; {\CalM} { I_3} 
 %\; dxdv dtdT 
 = - \nu  \int_0^{T_0} \int_0^T  \iint_{\T^3 \times \R^3}   |\nabla_v \; \mathcal{M} w| ^2. %dxdv dtdT.
$$


Since ${M}$ is uniformly bounded from above by $1$, we bound the term ${ I_{14}}$ as follows 
\begin{align*} 
& \quad \,
-3m \nu \int_0^{T_0}\int_0^{T} \iint_{\T^3 \times \R^3} { ( {\CalM} w)  }  {\CalM} { \frac{ w }{\langle v \rangle ^2} } 
\dx\dv\dt\dT
%\;  d\xi d\eta dtdT  
\\
&=3m \nu  \sum_{\eta} \int_0^{T_0}\int_0^{T}\int_{\R^3}  |M|^2 \mathcal{F}\left({ \frac{ w }{\langle v \rangle ^2} }\right) \mathcal{F}\left(w\right) \dxi \dt\dT 
\\
&\le 
3m \nu\int_0^{T_0}\int_0^{T}  \iint_{\T^3 \times \R^3} w^2 \dx\dv\dt\dT 
\leq 
3m \nu T_0 \int_0^{T_0}\iint_{\T^3 \times \R^3} w^2 \dx\dv\dt. 
\end{align*}
We carry out similar estimates for ${I_{13}}$.  To bound ${I_{12}}$ we first note that  
\begin{align*}
\frac{v \cdot \nabla w }{\langle v \rangle ^2} = \textrm{div} \left(\frac{v  w }{\langle v \rangle ^2} \right) - \frac{3w}{\langle v \rangle ^2} +  \frac{w|v|^2 }{\langle v \rangle ^2} =: J_1+J_2+J_3.
\end{align*}
The terms $J_2$ and $J_3$ are similar to $I_{14}$. For $J_1$ we have that 
\begin{align*}
& \quad \,
-2m\nu \int_0^{T_0}\int_0^{T} \iint_{\T^3 \times \R^3}  \CalM w \CalM \left( \textrm{div} \left(\frac{v  w }{\langle v \rangle ^2} \right)\right) \dx\dv\dt\dT 
\\
& = 2m\nu \int_0^{T_0}\int_0^{T} \iint_{\T^3 \times \R^3} \nabla \CalM w \cdot \CalM \left(\frac{v  w }{\langle v \rangle ^2} \right)  \dx\dv\dt\dT 
\\
& \leq \nu^2  \int_0^{T_0}\int_0^{T} \iint_{\T^3 \times \R^3} |\nabla \CalM w|^2 \dx\dv\dt\dT 
+ 4m^2 \int_0^{T_0}\int_0^{T} \iint_{\T^3 \times \R^3}  \left|\frac{v  w }{\langle v \rangle ^2} \right|^2  %\dx\dv\dt\dT 
\\
&\leq 
\nu^2  \int_0^{T_0}\int_0^{T} \iint_{\T^3 \times \R^3} |\nabla \CalM w|^2 \dx\dv\dt\dT 
+  4m^2 T_0 \int_0^{T_0} \iint_{\T^3 \times \R^3}  | w|^2  \dx\dv\dt.
\end{align*}
We now analyze $I_1=\textrm{div}_v \left( A[f_1] \nabla w - w \nabla a[f_1]\right)$.  Following the notation of Lemma \ref{lem:kernels}, we~write 
$$
\textrm{div}_v \left( A[f_1] \nabla w\right) = \textrm{div}_v \left( \vpran{A[f_1]\ast \psi_a} \nabla w\right) + \textrm{div}_v \left( (A[f_1]-A[f_1]\ast \psi_a )\nabla w\right).
$$
The matrix $ A[f_1]\ast \psi_a $ is nonnegative definite. In the term with the smooth diffusion we first integrate by parts and then commute $\mathcal{M}$ with $A[f_1]\ast \psi_a $. This yields  
\begin{align*}
& \quad \,
\int_0^{T_0}\int_0^{T} \iint_{\T^3 \times \R^3} \mathcal{M} w \; \mathcal{M} \; \textrm{div}_v \;\left( A[f_1]\ast \psi_a \nabla w \right)\dx\dv\dt\dT 
\\
&= -\int_0^{T_0}\int_0^{T} \iint_{\T^3 \times \R^3}  \nabla(\mathcal{M} w) \cdot  \; \mathcal{M} \; \;\left( A[f_1]\ast \psi_a \nabla w \right) \dx\dv\dt\dT 
\\
&= -\int_0^{T_0}\int_0^{T} \iint_{\T^3 \times \R^3}  \nabla(\mathcal{M} w) \cdot  \; [\mathcal{M}, A[f_1]\ast \psi_a] \nabla w \dx\dv\dt\dT
\\
& \quad \,
 -\int_0^{T_0}\int_0^{T} \iint_{\T^3 \times \R^3}  \nabla(\mathcal{M} w) \cdot  \; \left( A[f_1]\ast \psi_a \nabla (\mathcal{M}w) \right) \dx\dv\dt\dT.
\end{align*}
The last integral is nonnegative and can be discarded. For the first one, we use  Lemma \ref{lem:kernels} with $h:= A[f_1]\ast \psi_a$, Lemma \ref{lem:M-comm-x-v}, and Young's inequality to get 
\begin{align*}
& \quad \,
-\int_0^{T_0}\int_0^{T} \iint_{\T^3 \times \R^3}  \nabla(\mathcal{M} w) \cdot  \; [\mathcal{M}, A[f_1]\ast \psi_a] \nabla w \dx\dv\dt\dT 
\\
& \leq 
\frac{\nu}{20}\int_0^{T_0}\int_0^{T} \iint_{\T^3 \times \R^3}  |\nabla(\mathcal{M} w)|^2 %dxdvdtdT 
+ \frac{20}{\nu}\int_0^{T_0}\int_0^{T} \iint_{\T^3 \times \R^3} |\mathcal{T} \nabla  (\mathcal{M}w )|^2 
%+  |\mathcal{T}_2 \nabla  (\mathcal{M}w %)|^2\;dxdvdtdT 
\\
& \leq 
\frac{\nu}{20}\int_0^{T_0}\int_0^{T} \iint_{\T^3 \times \R^3}  
|\nabla(\mathcal{M} w)|^2 %dxdvdtdT 
+ \frac{20 c(a) }{\nu}\int_0^{T_0}\int_0^{T} \iint_{\T^3 \times \R^3} w^2 
%+  T  |\nabla  (\mathcal{M}w )|^2\;dxdvdtdT 
\\
&\le \left( \frac{\nu}{20} +  \frac{10 c(a)T_0 }{\nu}\right)  \int_0^{T_0}\int_0^{T} \iint_{\T^3 \times \R^3}  
|\nabla(\mathcal{M} w)|^2 %dxdvdtdT 
+ \frac{20 c(a) T_0}{\nu}\int_0^{T_0}\iint_{\T^3 \times \R^3} 
w^2.
%\;dxdvdt.
\end{align*}
We now look at the non-smooth part of the diffusion: 
\begin{align*}
& \quad \,
\int_0^{T_0}\int_0^{T} \iint_{\T^3 \times \R^3} 
\mathcal{M} w \; \mathcal{M}\textrm{div}_v \;\left( (A[f_1]-A[f_1]\ast \psi_a) \nabla w \right)\dx\dv\dt\dT
\\
&= - \int_0^{T_0}\int_0^{T} \iint_{\T^3 \times \R^3}  \nabla (\mathcal{M} w) \cdot \mathcal{M} \;\left( (A[f_1]-A[f_1]\ast \psi_a) \nabla w \right)\dx\dv\dt\dT
\\
&= - \int_0^{T_0}\int_0^{T} \iint_{\T^3 \times \R^3}  \nabla (\mathcal{M} w) \cdot \textrm{div}_v \mathcal{M} \;\left( (A[f_1]-A[f_1]\ast \psi_a)  w \right) \dx\dv\dt\dT
\\
& \quad \,
- \int_0^{T_0}\int_0^{T} \iint_{\T^3 \times \R^3}  \nabla (\mathcal{M} w) \cdot \mathcal{M} \;\left( (\nabla a[f_1]-\nabla a[f_1]\ast \psi_a)  w \right)\dx\dv\dt\dT =: B_1+B_2.
\end{align*}
In $B_1$, Young's inequality and (\ref{lapl}) yield
\begin{align*}
    |B_1| 
& \leq   
 \frac{\nu}{20} \int_0^{T_0}\int_0^{T} \iint_{\T^3 \times \R^3}  |\nabla(\mathcal{M} w)|^2 \dx\dv\dt\dT 
\\
& \quad \, 
    + \frac{20}{\nu} \sum_\eta \int_0^{T_0}\int_0^{T} \int_{\R^3}  |\xi|^2 M^2 |  \mathcal{F}({(A[f_1]- A[f_1]\ast \psi_a) w})|^2  \dxi \deta \dt\dT 
\\
& \leq  
   \frac{\nu}{20} \int_0^{T_0}\int_0^{T} \iint_{\T^3 \times \R^3}  |\nabla(\mathcal{M} w)|^2 \dx\dv\dt\dT 
\\
    & \quad \,
    +  \frac{40}{\nu \Eps \delta} \int_0^{T_0} \iint_{\T^3 \times \R^3}  |  (A[f_1]- A[f_1]\ast \psi_a) w|^2  \dx \dv \dt
\\
& \leq  
\frac{\nu}{20} \int_0^{T_0}\int_0^{T} \iint_{\T^3 \times \R^3}  |\nabla(\mathcal{M} w)|^2 \dx\dv\dt\dT 
\\
& \quad \, 
+  \frac{40}{\nu \Eps \delta}\|  (A[f_1]- A[f_1]\ast \psi_a)\|_{L^{\infty}_{t,x,v}}^2 \int_0^{T_0} \iint_{\T^3 \times \R^3}  w^2  \dx \dv \dt ,
\end{align*}
with the abuse of notation that $|  \mathcal{F}({(A[\cdot]- A[\cdot]\ast \psi_a) w})|^2$ stands for $\sum_{i=1}^3 |\mathcal{F}(A_i)|^2$ with $A_i$ the $i$-th column of the matrix $(A[\cdot]- A[\cdot]\ast \psi_a)w$.  Since $A[f_1]\ast \psi_a-A[f_1]$ converges to zero uniformly  in $(x,v,t)$  as $a \to 0$ for $f_1$ continuous and decaying fast enough for large velocity, we can chose  $a$ small enough such that 
$$
 \frac{40}{\nu \Eps \delta} \|  (A[f_1]- A[f_1]\ast \psi_a)\|_{L^{\infty}_{t,x,v}}^2 \le  \frac{1}{8}. 
$$
 Summarizing, 
 \begin{align*}
    |B_1|  \; \le \; &    \frac{\nu}{20} \int_0^{T_0}\int_0^{T} \iint_{\T^3 \times \R^3}  |\nabla(\mathcal{M} w)^2 \dx\dv\dt\dT  + \frac{1}{8} \int_0^{T_0} \iint_{\T^3 \times \R^3}  w^2  \dx \dv \dt.
 \end{align*}
In $B_2$ we apply Young's inequality and ${M}\le 1$ to get 
\begin{align*}
    |B_2| 
& \leq  
\frac{\nu}{20}\; \int_0^{T_0}\int_0^{T} \iint_{\T^3 \times \R^3}  |\nabla(\mathcal{M} w)|^2 %dxdvdtdT 
+  \frac{20}{\nu} \int_0^{T_0}\int_0^{T} \iint_{\T^3 \times \R^3} 
\abs{\nabla a[f_1]-\nabla a[f_1]\ast \psi_a}^2  w^2   %dxdvdtdT 
\\
&\leq  
\frac{\nu}{20} \int_0^{T_0}\int_0^{T} \iint_{\T^3 \times \R^3}  |\nabla(\mathcal{M} w)|^2 %dxdvdtdT 
+  \frac{20 \;T_0 }{\nu} \| \nabla a[f_1]-\nabla a[f_1]\ast \psi_a\|_{L^\infty_{t,x,v}}^2 \int_0^{T_0} \iint_{\T^3 \times \R^3}   w^2   %dxdvdt
\\
& \leq  
\frac{\nu}{20} \int_0^{T_0}\int_0^{T} \iint_{\T^3 \times \R^3}  
|\nabla(\mathcal{M} w)|^2 %dxdvdtdT 
+  \frac{20 c(m) T_0 }{\nu} \| g_1\|_{L^\infty_{t,x}(L^4_v)}^2 \int_0^{T_0} \iint_{\T^3 \times \R^3}   w^2.   %dxdvdt.
\end{align*}
The remaining part of $I_1$, namely $-\textrm{div}_v \left( w \nabla a[f_1]\right)$, can be estimated as $B_2$: 
\begin{align*}
& \quad \,
\int_0^{T_0}\int_0^{T} \iint_{\T^3 \times \R^3} \mathcal{M} w \; \mathcal{M}\textrm{div}_v \left( w \nabla a[f_1]\right) \dx\dv\dt\dT 
\\
& \leq  
\frac{\nu}{20}\; \int_0^{T_0}\int_0^{T} \iint_{\T^3 \times \R^3}  |\nabla(\mathcal{M} w)|^2 %dxdvdtdT  
+  \frac{20}{\nu} \int_0^{T_0}\int_0^{T} \iint_{\T^3 \times \R^3}   
|\nabla a[f_1]|^2  w^2   %dxdvdtdT 
\\
& \leq  
\frac{\nu}{20} \int_0^{T_0}\int_0^{T} \iint_{\T^3 \times \R^3}  
|\nabla(\mathcal{M} w)|^2 %dxdvdtdT 
+  \frac{20 c(m) T_0 }{\nu} \| g_1\|_{L^\infty_{t,x}(L^4_v)}^2 \int_0^{T_0} \iint_{\T^3 \times \R^3}   w^2.  % dxdvdt. 
\end{align*}
All the above yields 
\begin{align*}
    \int_0^{T_0}\int_0^{T} \iint_{\T^3 \times \R^3} \mathcal{M} w \; \mathcal{M}I_1 %\dx\dv\dt\dT 
& \leq  
\frac{3 \nu}{20}\; \int_0^{T_0}\int_0^{T} \iint_{\T^3 \times \R^3}  
|\nabla(\mathcal{M} w)|^2 \dx\dv\dt\dT 
\\
& \quad \,
 + \left( \frac{20 c(m) T_0 }{\nu} \| g_1\|_{L^\infty_{t,x}(L^4_v)}^2 + \frac{1}{8}\right) \int_0^{T_0} \iint_{\T^3 \times \R^3}   w^2   \dx\dv\dt. 
\end{align*}
We now look at $I_2 = \textrm{div}_v \left(A[f_1-f_2]\nabla g_2 \right)- \textrm{div}_v (g_2 \nabla a[f_1-f_2])$; for the second part we have 
\begin{align}
& \quad \, 
\int_0^{T_0} \int_0^{T} \iint_{\T^3 \times \R^3}  \mathcal{M} w \;\textrm{div}_v ( \mathcal{M} g_2 \nabla a[f_1-f_2])q \dx\dv\dt\dT \nn
\\
&\leq \frac{\nu}{20}\int_0^{T_0} \int_0^{T} \iint_{\T^3 \times \R^3}  |\nabla(\mathcal{M} w)|^2 %\; dxdvdtdT 
+ \frac{20}{\nu} \sum_\eta \int_0^{T_0} \int_0^{T} \int_{\R^3}  {M}^2 \abs{\mathcal{F} (g_2 \nabla a[f_1-f_2]}^2 \nn %\;d\xi d\eta dtdT \nonumber 
\\
&\leq 
\frac{\nu}{20}\int_0^{T_0} \int_0^{T} \iint_{\T^3 \times \R^3}  |\nabla(\mathcal{M} w)|^2 %\; dxdvdtdT 
+  \frac{20}{\nu}\int_0^{T_0} \int_0^{T} \iint_{\T^3 \times \R^3}   g_2^2 |\nabla a[f_1-f_2]|^2 %\;dx dv dtdT  
\nonumber 
\\
&\leq \frac{\nu}{20}\int_0^{T_0} \int_0^{T} \iint_{\T^3 \times \R^3}  |\nabla(\mathcal{M} w)|^2 %\; dxdvdtdT 
+ \frac{20}{\nu}\|g_2\|^2_{L_{x,t}^\infty (L^3_v)}
 \int_0^{T_0} \int_0^{T}\int_{\T^3} \left(  \int_{\mathbb{R}^3 }|\nabla a[f_1-f_2] |^6 \;dv \right)^{1/3} %\;dxdtdT    
\nonumber 
\\
&\leq 
\frac{\nu}{20}\int_0^{T_0} \int_0^{T} \iint_{\T^3 \times \R^3}  |\nabla(\mathcal{M} w)|^2 %\; dxdvdtdT 
+ \frac{20 c(m)}{\nu}\|g_2\|^2_{L_{x,t}^\infty (L^3_v)} 
 \int_0^{T_0} \int_0^{T}\iint_{\T^3 \times \R^3} |f_1-f_2 |^2 %\;dvdxdtdT  
 \nonumber 
\\
&\leq 
\frac{\nu}{20} \int_0^{T_0}\int_0^{T} \iint_{\T^3 \times \R^3}  
|\nabla(\mathcal{M} w)|^2 %\; dxdvdtdT     
+ \frac{20 c(m)  \|g_2\|^2_{L_{x,t}^\infty (L^3_v)}}{\nu}T_0  \int_0^{T_0}\iint_{\T^3 \times \R^3}  | w|^2, \label{I_2_partII} %\;dvdxdt,  
%\nonumber
\end{align}
using (\ref{nablaa_L6}) and ${M}\le 1$. For the diffusion term, we apply the same decomposition between smooth and less smooth parts as in $I_1$, this time for $\nabla g_2$: 
$$
\textrm{div}_v \left( A[f_1-f_2] \nabla g_2\right) = \textrm{div}_v \left( A[f_1-f_2] \nabla g_2 \ast \psi_a \right) + \textrm{div}_v \left( (A[f_1-f_2] (\nabla g_2 - \nabla g_2 \ast \psi_a )\right).
$$
We estimate the smooth part as 
\begin{align*}
& \quad \,
   \int_0^{T_0} \int_0^{T} \iint_{\T^3 \times \R^3}  \mathcal{M} w \; \mathcal{M} \textrm{div}_v \left( A[f_1-f_2] \nabla g_2 \ast \psi_a \right) \dx\dv\dt\dT  
\\
& = - \int_0^{T_0} \int_0^{T} \iint_{\T^3 \times \R^3}  \nabla(\mathcal{M} w) \cdot  \mathcal{M} \left( A[f_1-f_2] g_2 \ast \nabla \psi_a \right) \dx\dv\dt\dT 
\\
& \leq  
\frac{\nu}{20}\int_0^{T_0} \int_0^{T} \iint_{\T^3 \times \R^3}  |\nabla(\mathcal{M} w)|^2 %\; dxdvdtdT 
+  \frac{20}{\nu} \sum_\eta \int_0^{T_0} \int_0^{T} \int_{\R^3} M^2 | \mathcal{F}( A[f_1-f_2] g_2 \ast \nabla \psi_a)|^2 
%d\eta d\xi dtdT
\\
&\leq 
\frac{\nu}{20}\int_0^{T_0} \int_0^{T} \iint_{\T^3 \times \R^3}  |\nabla(\mathcal{M} w)|^2 %\; dxdvdtdT 
+  \frac{20}{\nu} \int_0^{T_0} \int_0^{T} \iint_{\T^3 \times \R^3}  
|  A[f_1-f_2] g_2 \ast \nabla \psi_a|^2,
%dxdv dtdT,
\end{align*}
using once more that $M\le 1$. In the last integral, we use the smoothness of $g_2 \ast \nabla \psi_a$ to obtain 
\begin{align*}
   \int_0^{T_0} \int_0^{T}& \iint_{\T^3 \times \R^3}  |  A[f_1-f_2] g_2 \ast \nabla \psi_a|^2 \dx\dv \dt\dT
\\
& \leq   \int_0^{T_0} \int_0^{T} \int_{\T^3} \left(\| A[f_1-f_2] \|^2_{L^\infty_v} \int_{\R^3}  |g_2 \ast \nabla \psi_a|^2 \dv\right) \dx\dt\dT
\\
& \leq  
c(a) c(m) \int_0^{T_0} \int_0^{T} \int_{\T^3} \| w \|^2_{L^2_v} \|g_2\|^2_{L^1_v} \dx\dt\dT
\\
& \le  
c(a) c(m) \|g_2\|^2_{L^\infty_{t,x}L^1_v}  \int_0^{T_0} \int_0^{T} \iint_{\T^3 \times \R^3} w^2 \dx\dv\dt\dT
\\
& \le  
c(a) c(m)  T_0 \|g_2\|^2_{L^\infty_{t,x}L^1_v}  \int_0^{T_0} \iint_{\T^3 \times \R^3} w^2 \dx\dv\dt. 
\end{align*}
Hence we get 
\begin{align*}
& \quad \,
    \int_0^{T_0} \int_0^{T} \iint_{\T^3 \times \R^3}   \mathcal{M} w \; \mathcal{M} \textrm{div}_v \left( A[f_1-f_2] \nabla g_2 \ast \psi_a \right) \; \dx\dv\dt\dT 
\\
&\leq  
\frac{\nu}{20}\int_0^{T_0} \int_0^{T} \iint_{\T^3 \times \R^3}  
|\nabla(\mathcal{M} w)|^2 %\; dxdvdtdT 
+  \frac{20}{\nu}  c(a) \, T_0 \, \|g_2\|^2_{L^\infty_{t,x}L^1_v}  \int_0^{T_0}  \iint_{\T^3 \times \R^3} w^2. %dxdvdt.
\end{align*}
The less smooth part can be estimated as 
\begin{align*}
   \int_0^{T_0} \int_0^{T}& \iint_{\T^3 \times \R^3}    \mathcal{M} w \; \mathcal{M} \textrm{div}_v \left( A[f_1-f_2](\nabla g_2 - \nabla g_2 \ast \psi_a )\right) \; \dx\dv\dt\dT  
\\
   =& - \int_0^{T_0} \int_0^{T} \iint_{\T^3 \times \R^3}  \nabla(\mathcal{M} w) \cdot  \mathcal{M} \left( A[f_1-f_2] (\nabla g_2 - \nabla g_2 \ast \psi_a) \right) \dx\dv\dt\dT 
\\
&= - \int_0^{T_0} \int_0^{T} \iint_{\T^3 \times \R^3}  \nabla(\mathcal{M} w) \cdot  \textrm{div}_v ( \mathcal{M} \left( A[f_1-f_2] ( g_2 -g_2 \ast \psi_a) \right) \dx\dv\dt\dT 
\\
& +  \int_0^{T_0} \int_0^{T} \iint_{\T^3 \times \R^3}  \nabla(\mathcal{M} w) \cdot   \mathcal{M} \vpran{ \nabla a[f_1-f_2] \; ( g_2 -g_2 \ast \psi_a)}  \dx\dv\dt\dT. 
\end{align*}
The second term can be estimated as in (\ref{I_2_partII}). For the first, we apply Young's inequality and (\ref{int_xi_M2}) to get 
\begin{align*}
    - \int_0^{T_0}& \int_0^{T} \iint_{\T^3 \times \R^3}  \nabla(\mathcal{M} w) \cdot  \textrm{div}_v ( \mathcal{M} \left( A[f_1-f_2] ( g_2 -g_2 \ast \psi_a) \right) \dx\dv\dt\dT 
\\
    \le  \; & \frac{\nu}{20}\int_0^{T_0} \int_0^{T} \iint_{\T^3 \times \R^3}  |\nabla(\mathcal{M} w)|^2 \dx\dv\dt\dT  
\\
   &  + \frac{40}{\nu \Eps \delta } \int_0^{T_0} \iint_{\T^3 \times \R^3} | A[f_1-f_2] ( g_2 -g_2 \ast \psi_a)|^2 \dx\dv\dt 
\\
 \le  \; & \frac{\nu}{20}\int_0^{T_0} \int_0^{T} \iint_{\T^3 \times \R^3}  |\nabla(\mathcal{M} w)|^2 \dx\dv\dt\dT  
\\
   &  +  \frac{20 c(m)}{\nu \Eps \delta} \|  (g_2 -g_2 \ast \psi_a)\|^2_{L^\infty_{t,x}(L^2_v)} \int_0^{T_0} \iint_{\T^3 \times \R^3} w^2 \dx\dv\dt    
\\
    \le  \; & \frac{\nu}{20}\int_0^{T_0} \int_0^{T} \iint_{\T^3 \times \R^3}  |\nabla(\mathcal{M} w)|^2 %\dx\dv\dt\dT 
+  \frac{1}{8 } \int_0^{T_0} \iint_{\T^3 \times \R^3} w^2, %\;dxdvdt,
\end{align*}
for $a$ small enough, thanks to the uniform convergence of $g_2 -g_2 \ast \psi_a$ to zero as $a \to 0$.
Summarizing, we have 
\begin{align*}
    \int_0^{T_0}\int_0^{T} \iint_{\T^3 \times \R^3} \mathcal{M} w \; \mathcal{M}I_2 \dx\dv\dt\dT 
\le & \;  
\frac{ \nu}{5}\; \int_0^{T_0}\int_0^{T} \iint_{\T^3 \times \R^3}  |\nabla(\mathcal{M} w)|^2 \dx\dv\dt\dT 
\\
    &+ \left(c \, T_0+ \frac{1}{8}\right) \int_0^{T_0} \iint_{\T^3 \times \R^3}   w^2  \dx\dv\dt,
\end{align*}
where $c$ depends on $\nu$, some $L^\infty_{t,x}(L^2_v)$ norms of $g_i$ and other universal constants. 

%%%%%%%%%%%%%%%%%%%%
%%%%%%%%%%%%%%%%%%%%
%%%%%%%%%%%%%%%%%%%%%%

The estimates of ${I_4}$ to ${I_{11}}$ are similar, and we only sketch them briefly. For ${I_4}$ we have 
\begin{align*}
   \int_0^{T_0}\int_0^{T} \iint_{\T^3 \times \R^3}  \CalM w \CalM I_4 %dxdvdtdT  
= m  \int_0^{T_0}\int_0^{T} \iint_{\T^3 \times \R^3}  \nabla (\CalM w ) \CalM\left(\frac{A[f_1]v  w }{\langle v \rangle ^2} \right) %dxdvdtdT 
\\
\le  
\frac{\nu}{10}  \int_0^{T_0}\int_0^{T} \iint_{\T^3 \times \R^3}  |\nabla (\CalM w )| ^2 %dxdvdtdT 
+ \frac{10 \, c(m) m^2 \, T_0 \, \|g_1\|^2_{L^\infty_{t,x}L^2_v}}{\nu}  \int_0^{T_0} \iint_{\T^3 \times \R^3}  | {w} |^2, %dxdv dt,
\end{align*}
and similar estimates hold for the term with $I_5$:  % FROM HERE! 
\begin{align*}
& \quad \,
\int_0^{T_0}\int_0^{T} \iint_{\T^3 \times \R^3} 
 {({\CalM}w)} {\CalM} I_5 \dx\dv\dt\dT
\\ 
& \le \;  \frac{\nu}{10}  \int_0^{T_0}\int_0^{T} \iint_{\T^3 \times \R^3}  
|\nabla (\CalM w )| ^2 %\;dxdvdtdT 
+ \; \frac{10 \, m^2}{\nu}  \int_0^{T_0}  \int_0^T \iint_{\T^3 \times \R^3} 
|A[f_1-f_2]  g_2|^2 %\; dxdv dt dT 
\\
&\le \;  \frac{\nu}{10}  \int_0^{T_0}\int_0^{T} \iint_{\T^3 \times \R^3}  |\nabla (\CalM w )| ^2 %\;dxdvdtdT 
+ \; \frac{10 \, c(m)m^2 \, T_0 \, \|g_2\|^2_{L^\infty_{t,x}L^2_v}}{\nu}  \int_0^{T_0} \iint_{\T^3 \times \R^3}  |w|^2. %dxdvdt.
\end{align*}
For $I_{11}$ we also use the fact that $ {M} \le 1$ to get 
\begin{align*}
& \quad \,
\int_0^{T_0}\int_0^{T} \iint_{\T^3 \times \R^3} { ( {\CalM} w)  } \CalM {I_{11}} \dx\dv\dt\dT  
\\
&\le 
\int_0^{T_0}\int_0^{T} \iint_{\T^3 \times \R^3} |w| ^2 \dx\dv\dt\dT 
+ \int_0^{T_0}\int_0^{T} \iint_{\T^3 \times \R^3} \left| I_{11}\right|^2 \dx\dv\dt\dT
\\
&\le 
T_0 \int_0^{T_0} \iint_{\T^3 \times \R^3}  | w| ^2 \dx\dv\dt  
+  T_0 \int_0^{T_0} \iint_{\T^3 \times \R^3}  g_2^2 | \nabla a[f_1-f_2]|^2 \dx\dv\dt 
\\
&\le   
T_0 \int_0^{T_0} \iint_{\T^3 \times \R^3} |w| ^2 %dxdv dt 
+  T_0 \int_0^{T_0} \int_{\mathbb{R}^3} \|g_2\|_{L_v^3}^2\left( \int_{\mathbb{R}^3} | \nabla a[f_1-f_2]|^6 \dv \right)^{1/3} \dx\dt
\\
&\le   
T_0 \vpran{1+ c(m) \|g_2\|^2_{L^\infty_{x,t}L_v^3}} \int_0^{T} \iint_{\T^3 \times \R^3} w^2 \dx\dv\dt.
\end{align*}
Similarly, for the term with $I_9$ we have
\begin{align*}
\int_0^{T_0}\int_0^{T} & \iint_{\T^3 \times \R^3} {({\CalM} w)  } \CalM {I_{9}} \dx\dv\dt\dT    
\\
&\le 
T_0 \int_0^{T_0} \iint_{\T^3 \times \R^3} |w| ^2 \dx\dv\dt 
+  T_0 \int_0^{T_0} \iint_{\T^3 \times \R^3}  g_2^2 | A[f_1-f_2]|^2 \dx\dv\dt 
\\
&\le 
T_0 \int_0^{T_0} \iint_{\T^3 \times \R^3} |w| ^2 \dx\dv\dt 
+  T_0 c(m) \|g_2\|^2_{L^\infty_{x,t}L_v^2} \int_0^{T_0} \iint_{\T^3 \times \R^3} w^2 \dx\dv\dt
\\
&\le T_0 \vpran{1+c(m) \|g_2\|^2_{L^\infty_{t,x}L_v^2}} \int_0^{T_0} \iint_{\T^3 \times \R^3} w^2 \dx\dv\dt,
\end{align*}
using (\ref{A_sup}) to estimate $A[f_1-f_2]$. 
We rewrite $I_6$ and $I_7$ as 
\begin{align*}
-\left  \langle A[f_1]   \nabla w ,  \frac{v}{ \langle v \rangle ^2}\right \rangle = &\;   {{w\; \textrm{div}_v\left( A[f_1] \frac{v}{ \langle v \rangle ^2}   \right)}}
- \textrm{div} \left(  w A[f_1]    \frac{v}{ \langle v \rangle ^2} \right) 
\\
=&:\; w I_{6,1}- I_{6,2},
\\
 -  \left  \langle A[f_1-f_2]   \nabla g_2 ,  \frac{v}{ \langle v \rangle ^2}\right  \rangle = &\;   {{g_2\; \textrm{div}_v\left( A[f_1-f_2] \frac{v}{ \langle v \rangle ^2}   \right)}}
- \textrm{div} \left(  g_2 A[f_1-f_2]    \frac{v}{ \langle v \rangle ^2} \right) 
\\
=\; g_2 \bigg(\nabla a [f_1-f_2] & \frac{v}{ \langle v \rangle ^2} + Tr(A[f_1-f_2] \nabla \frac{v}{ \langle v \rangle ^2})\bigg) - \textrm{div} \left(  g_2 A[f_1-f_2]    \frac{v}{ \langle v \rangle ^2} \right) 
\\
=: \;  g_2 I_{7,1} + g_2 & \, I_{7,2} -  I_{7,3}.
\end{align*}
The term $I_{6,2}$ is similar to $I_{4}$.  For $I_{6,1}$ we use the bound $ {M} \le 1$ and get
\begin{align*}
\int_0^{T_0}\int_0^{T} & \iint_{\T^3 \times \R^3} \abs{\CalM w} |{\CalM}  \vpran{w I_{6,1}}| %dxdv dtdT  
\leq 
\left(\int_0^{T_0}\int_0^{T} \iint_{\T^3 \times \R^3} | w| ^2 \right)^{1/2} 
\left( \int_0^{T_0}\int_0^{T} \iint_{\T^3 \times \R^3}  |{w I_{6,1}}|^2 %dxdv
\right)^{1/2}   
\\
& \le  
\|I_{6,1}\|_{L^\infty_{t,x,v}} \int_0^{T_0}\int_0^{T} \iint_{\T^3 \times \R^3}  |w|^2 \dx\dv\dt\dT  
\\
&\le \|A[f_1] + \nabla a[f_1]\|_{L^\infty_{t,x,v}} T_0   \int_0^{T_0}\iint_{\T^3 \times \R^3}  |w| ^2 \dx\dv\dt
\\
&\le c(m) \|g_1\|_{L^\infty_{t,x}(L_v^1\cap L_v^4)} T_0 \int_0^{T_0}\iint_{\T^3 \times \R^3}  |w| ^2 \dx\dv\dt. 
\end{align*}
The term $g_2 I_{7,1}$ can be estimated as $I_{11}$, $g_2 I_{7,2}$ as $I_9$, and $I_{7,3}$ as $I_5$.  
Finally, $I_8$ and $I_{10}$ can be estimated as $w I_{6,1}$. Summarizing all the estimates for $I_1,...,I_{14}$ we get 
\begin{align*}
\left(\frac{1}{4}-cT_0\right) \int_0^{T_0} \iint_{\T^3 \times \R^3} w^2 \dx\dv\dt 
\le 
-\left(\frac{\nu}{10}-\delta\left(\frac{1}{2}+\Eps\right)\right) \int_0^{T_0}\int_0^{T} \iint_{\T^3 \times \R^3} |\nabla(\CalM w)|^2, %\;dxdvdtdT  
\end{align*}
for $c$ a constant depending on $m$, $\nu$, $\Eps$, $\delta$ and $\|g_i\|_{L^\infty_{t,x}(L^1_v\cap L^4_v)}$. We pick $\delta$ small enough such that $\left(\frac{\nu}{10}-\delta\left(\frac{1}{2}+\Eps\right)\right)\le 0$ and then $T_0$ small enough so that $\frac{1}{4}-cT_0\ge 0$ to conclude that $w = 0$ for all times $t\in [0,T_0]$. 
\end{proof}

\end{document}
